\documentclass[10pt,a4paper]{amsart}
\usepackage{amsmath,amssymb,amsthm,graphicx}
% Retained minimal preamble for the Stacks Project chapter more-algebra.tex (More on Algebra).
% The chapter file pinned-original-source.tex carries only \input{preamble} on its first line, so no
% preamble text lives inside the chapter itself.  The Stacks Project uses one single shared preamble for
% every chapter; the snapshot stages it identically as sources/stacks-chow/preamble.tex and
% sources/stacks-etale-cohomology/preamble.tex.  The block below is copied mechanically, line for line,
% from sources/stacks-chow/preamble.tex lines 14-18 (the Xy-pic setup) and 154-217 (the native theorem
% declarations together with the mathematical macros).  Nothing was added, removed or redefined; the
% wrapper class and the amsmath/amssymb/amsthm/graphicx base packages are supplied by the assembler.
% Chapter: more-algebra.tex -- More on Algebra.
% For commutative diagrams we use Xy-pic
\usepackage[all]{xy}

% We use 2cell for 2-commutative diagrams.
\xyoption{2cell}

% Theorem environments.
%
\theoremstyle{plain}
\newtheorem{theorem}[subsection]{Theorem}
\newtheorem{proposition}[subsection]{Proposition}
\newtheorem{lemma}[subsection]{Lemma}

\theoremstyle{definition}
\newtheorem{definition}[subsection]{Definition}
\newtheorem{example}[subsection]{Example}
\newtheorem{exercise}[subsection]{Exercise}
\newtheorem{situation}[subsection]{Situation}

\theoremstyle{remark}
\newtheorem{remark}[subsection]{Remark}
\newtheorem{remarks}[subsection]{Remarks}

\numberwithin{equation}{subsection}

% Macros
%
\def\lim{\mathop{\mathrm{lim}}\nolimits}
\def\colim{\mathop{\mathrm{colim}}\nolimits}
\def\Spec{\mathop{\mathrm{Spec}}}
\def\Hom{\mathop{\mathrm{Hom}}\nolimits}
\def\Ext{\mathop{\mathrm{Ext}}\nolimits}
\def\SheafHom{\mathop{\mathcal{H}\!\mathit{om}}\nolimits}
\def\SheafExt{\mathop{\mathcal{E}\!\mathit{xt}}\nolimits}
\def\Sch{\mathit{Sch}}
\def\Mor{\mathop{\mathrm{Mor}}\nolimits}
\def\Ob{\mathop{\mathrm{Ob}}\nolimits}
\def\Sh{\mathop{\mathit{Sh}}\nolimits}
\def\NL{\mathop{N\!L}\nolimits}
\def\CH{\mathop{\mathrm{CH}}\nolimits}
\def\proetale{{pro\text{-}\acute{e}tale}}
\def\etale{{\acute{e}tale}}
\def\QCoh{\mathit{QCoh}}
\def\Ker{\mathop{\mathrm{Ker}}}
\def\Im{\mathop{\mathrm{Im}}}
\def\Coker{\mathop{\mathrm{Coker}}}
\def\Coim{\mathop{\mathrm{Coim}}}

% Boxtimes
%
\DeclareMathSymbol{\boxtimes}{\mathbin}{AMSa}{"02}

%
% Macros for moduli stacks/spaces
%
\def\QCohstack{\mathcal{QC}\!\mathit{oh}}
\def\Cohstack{\mathcal{C}\!\mathit{oh}}
\def\Spacesstack{\mathcal{S}\!\mathit{paces}}
\def\Quotfunctor{\mathrm{Quot}}
\def\Hilbfunctor{\mathrm{Hilb}}
\def\Curvesstack{\mathcal{C}\!\mathit{urves}}
\def\Polarizedstack{\mathcal{P}\!\mathit{olarized}}
\def\Complexesstack{\mathcal{C}\!\mathit{omplexes}}
% \Pic is the operator that assigns to X its picard group, usage \Pic(X)
% \Picardstack_{X/B} denotes the Picard stack of X over B
% \Picardfunctor_{X/B} denotes the Picard functor of X over B
\def\Pic{\mathop{\mathrm{Pic}}\nolimits}
\def\Picardstack{\mathcal{P}\!\mathit{ic}}
\def\Picardfunctor{\mathrm{Pic}}
\def\Deformationcategory{\mathcal{D}\!\mathit{ef}}
\usepackage{cleveref}
\crefname{originallemma}{original lemma}{original lemmas}
\Crefname{originallemma}{Original lemma}{Original lemmas}
\makeatletter
\newlabel{lemma-derived-tor-quasi-isomorphism}{{064K}{source}{}{sourcepointer-1}{}}
\newlabel{lemma-derived-tor-quasi-isomorphism@cref}{{[originallemma][0][]064K}{[1][0][]source}{}{}{}}
\makeatother
\crefname{originallemma}{original lemma}{original lemmas}
\Crefname{originallemma}{Original lemma}{Original lemmas}
\makeatletter
\newlabel{lemma-cone-relatively-pseudo-coherent}{{0674}{source}{}{sourcepointer-2}{}}
\newlabel{lemma-cone-relatively-pseudo-coherent@cref}{{[originallemma][0][]0674}{[1][0][]source}{}{}{}}
\makeatother
\crefname{originallemma}{original lemma}{original lemmas}
\Crefname{originallemma}{Original lemma}{Original lemmas}
\makeatletter
\newlabel{lemma-base-change-relative-pseudo-coherent}{{067A}{source}{}{sourcepointer-3}{}}
\newlabel{lemma-base-change-relative-pseudo-coherent@cref}{{[originallemma][0][]067A}{[1][0][]source}{}{}{}}
\makeatother
\title{Relative pseudo-coherence is preserved by derived base change along a relatively pseudo-coherent algebra}
\author{The Stacks Project Authors}
\date{stacks-project:more-algebra.tex:local-snapshot-20261012; GNU FDL 1.2 or later}
\begin{document}
\maketitle
\setcounter{subsection}{1}
\noindent\textit{Editorial note.} All delivered mathematical text is the author's own, copied from the pinned Stacks Project chapter file more-algebra.tex; no statement or proof step was rewritten, reordered or removed. The counted statement is the labeled lemma environment at source lines 21567-21576 (canonical tag 067B) and the counted proof is the proof environment at source lines 21578-21619, which immediately follows it and is its own proof.\par\medskip
\noindent\textit{Editorial note.} The proof refers three times to earlier lemmas of the same chapter, namely the quasi-isomorphism for derived Tor, the cone criterion for relative pseudo-coherence and the base change lemma; all three are carried as navigation stubs with their roles recorded in the reasons.\par\medskip
\noindent\textit{Editorial note.} The author writes the polynomial variables as y\_1, ..., y\_m in the presentations and as y\_1, ..., y\_n once in the resolution, and the derived category statement is written with the symbol for the total complex; both are copied exactly as printed.\par\medskip
\noindent\textit{Editorial note.} The chapter file pinned-original-source.tex carries only an input of the shared preamble on its first line, so no preamble text lives inside the chapter. The retained file sources/stacks-more-algebra/retained-preamble.tex is a mechanical, line-for-line copy of the project's single shared preamble as staged in the snapshot (Xy-pic setup plus the native theorem declarations and mathematical macros). Nothing was added, removed or redefined, and the excerpt serves only as preamble, never as a span coordinate.\par\medskip
\noindent\textit{Editorial note.} The source URL is the source-level Stacks Project URL, as required when more than one card is produced from one source; the per-claim tag is carried by claim\_key and by this note. The counted statement and proof cite no bibliography entry, so this card needs no bibliography.\par\medskip
\noindent\textit{Allowed standard prerequisites (editorial).}\begin{itemize}
\item A relatively pseudo-coherent algebra admits a resolution by a complex of finite free modules over a polynomial algebra, and a bounded above complex may be replaced by one of flat modules when testing relative pseudo-coherence. used at proof lines 21579-21585. Produces the resolution E and the flat representative of K that the displayed computation uses.
\item Tensoring with a bounded above complex of flat modules computes the derived tensor product, and replacing the second factor by a quasi-isomorphic complex of flat modules does not change the total complex up to quasi-isomorphism. used at proof lines 21586-21601. Justifies the chain of equalities and the quasi-isomorphism in the display that rewrites the derived base change as a total complex over R.
\item A total complex filtered by subcomplexes whose successive quotients are m-pseudo-coherent is m-pseudo-coherent, and the cone criterion propagates relative pseudo-coherence along such filtrations. used at proof lines 21605-21612. Reduces the assertion for the total complex to the assertions for the single filtration quotients.
\end{itemize}\medskip
\ifdefined\hypertarget\hypertarget{sourcepointer-1}{}\fi
\noindent\textit{Source reference.} Original lemma 064K: sources/stacks-more-algebra/pinned-original-source.tex, source lines 14284--14284. The reference supplies the quasi-isomorphism that replaces the second factor by the finite free resolution E in the total complex, used at proof lines 21599-21601; the display and the filtration argument that follow are inside the counted span.\par\medskip
\ifdefined\hypertarget\hypertarget{sourcepointer-2}{}\fi
\noindent\textit{Source reference.} Original lemma 0674: sources/stacks-more-algebra/pinned-original-source.tex, source lines 21393--21393. The reference is the cone criterion invoked at proof lines 21610-21612 to pass from the filtration quotients to the total complex; the filtration and the reduction to a single quotient are inside the counted span.\par\medskip
\ifdefined\hypertarget\hypertarget{sourcepointer-3}{}\fi
\noindent\textit{Source reference.} Original lemma 067A: sources/stacks-more-algebra/pinned-original-source.tex, source lines 21543--21543. The reference is the base change statement for a polynomial algebra extension, invoked at proof lines 21616-21618 as the last remaining case; the reduction of the general case to that one is inside the counted span.\par\medskip
\begin{lemma}
\label{lemma-pull-relative-pseudo-coherent}
Let $R \to A \to B$ be finite type ring maps.
Let $m \in \mathbf{Z}$.
Let $K^\bullet$ be a complex of $A$-modules.
Assume $B$ as a $B$-module is pseudo-coherent relative to $A$.
If $K^\bullet$ is $m$-pseudo-coherent (resp.\ pseudo-coherent)
relative to $R$, then $K^\bullet \otimes_A^{\mathbf{L}} B$ is
$m$-pseudo-coherent (resp.\ pseudo-coherent) relative to $R$.
\end{lemma}

\begin{proof}
Choose a surjection $A[y_1, \ldots, y_m] \to B$.
Choose a surjection $R[x_1, \ldots, x_n] \to A$.
Combined we get a surjection $R[x_1, \ldots, x_n, y_1, \ldots y_m] \to B$.
Choose a resolution $E^\bullet \to B$ of $B$ by a complex of
finite free $A[y_1, \ldots, y_n]$-modules (which is possible
by our assumption on the ring map $A \to B$). We may assume
that $K^\bullet$ is a bounded above complex of flat $A$-modules. Then
\begin{align*}
K^\bullet \otimes_A^{\mathbf{L}} B & =
\text{Tot}(K^\bullet \otimes_A B[0]) \\
& = \text{Tot}(K^\bullet \otimes_A A[y_1, \ldots, y_m]
\otimes_{A[y_1, \ldots, y_m]} B[0]) \\
& \cong
\text{Tot}\left(
(K^\bullet \otimes_A A[y_1, \ldots, y_m])
\otimes_{A[y_1, \ldots, y_m]} E^\bullet
\right) \\
& =
\text{Tot}(K^\bullet \otimes_A E^\bullet)
\end{align*}
in $D(A[y_1, \ldots, y_m])$. The quasi-isomorphism $\cong$ comes from
an application of
Lemma \ref{lemma-derived-tor-quasi-isomorphism}.
Thus we have to show that
$\text{Tot}(K^\bullet \otimes_A E^\bullet)$ is $m$-pseudo-coherent
as a complex of $R[x_1, \ldots, x_n, y_1, \ldots y_m]$-modules.
Note that $\text{Tot}(K^\bullet \otimes_A E^\bullet)$ has a filtration by
subcomplexes with successive quotients the complexes
$K^\bullet \otimes_A E^i[-i]$. Note that for $i \ll 0$ the
complexes $K^\bullet \otimes_A E^i[-i]$ have zero cohomology
in degrees $\leq m$ and hence are $m$-pseudo-coherent (over any ring).
Hence, applying
Lemma \ref{lemma-cone-relatively-pseudo-coherent}
and induction, it suffices to show that $K^\bullet \otimes_A E^i[-i]$ is
pseudo-coherent relative to $R$ for all $i$. Note that $E^i = 0$ for
$i > 0$. Since also $E^i$ is finite free this
reduces to proving that $K^\bullet \otimes_A A[y_1, \ldots, y_m]$ is
$m$-pseudo-coherent relative to $R$ which follows from
Lemma \ref{lemma-base-change-relative-pseudo-coherent}
for instance.
\end{proof}

\end{document}
